\documentclass[a4paper]{article}

\usepackage{graphicx}
\usepackage{amsmath,amssymb}
\usepackage{minted}
\usepackage{multirow}
\usepackage{float}
\usepackage{algorithm}
\usepackage{algpseudocode}
\usepackage{todonotes}
\usepackage{forest}
\usepackage{xstring}
\usepackage{xcolor}
\usepackage[most]{tcolorbox}
\usepackage{xparse}
\usepackage{natbib}
\usepackage{adjustbox}

\usetikzlibrary{graphs,graphdrawing}
\usegdlibrary{layered}

% for external graphviz
\input{/home/donkarlo/Dropbox/repo/nd_language_ai_project/src/nd_language_ai/formal/computer/programming/tex/latex/code_clips/external_graphviz.tex}

% for tags
\input{/home/donkarlo/Dropbox/repo/nd_language_ai_project/src/nd_language_ai/formal/computer/programming/tex/latex/part/control_sequence/kind/semtag/semtag.tex}

% for links
\input{/home/donkarlo/Dropbox/repo/nd_language_ai_project/src/nd_language_ai/formal/computer/programming/tex/latex/code_clips/seehere.tex}

% for nested lists and sections
\input{/home/donkarlo/Dropbox/repo/nd_language_ai_project/src/nd_language_ai/formal/computer/programming/tex/latex/code_clips/deep_structure.tex}

\title{Self-other distinction}
\date{}
\author{Mohammad Rahmani}

\begin{document}
   \maketitle

   \section{Definition}
         \cite{decety-2003-shared-representations-between-self-and-other,tsakiris-2007-on-agency-and-body-ownership} describe self--other distinction as the ability to maintain a distinction between representations related to oneself and representations related to other agents, despite partially shared perceptual and motor representations. In the bodily domain, this distinction is closely related to the sense of body ownership and the sense of agency over actions.

   \section{Approaches}

         \subsubsection{Visual--proprioceptive and action congruence}
            \cite{van-den-bos-2002-sense-of-body-and-sense-of-action-both-contribute-to-self-recognition} showed that self-recognition depends on congruence between visual information and signals about body position, as well as on information about one's own actions. Participants distinguished their own hand from another person's hand much more accurately when the two hands performed different movements. The results support sensorimotor congruence as a mechanism for separating self-generated bodily events from those generated by others.

         \subsubsection{Efference-based self-recognition}
            \cite{tsakiris-2005-a-specific-role-for-efferent-information-in-self-recognition} investigated whether efferent motor information contributes specifically to self-recognition. The study found that self-generated actions provide information that improves discrimination between one's own body and another person's body, beyond visual feedback alone. This supports the use of motor commands or their internal copies as cues for attributing observed actions to the self.

         \subsubsection{Sense of agency and body ownership}
            \cite{kalckert-2012-moving-a-rubber-hand-that-feels-like-your-own} used a moving rubber-hand illusion to separate the sense of agency from the sense of body ownership. Synchronous visuomotor feedback could induce both the feeling that the artificial hand belonged to the participant and the feeling of controlling its movement, while experimental manipulations dissociated these two experiences. The findings indicate that self-attribution of a limb and self-attribution of an action are related but distinct processes.

         \subsubsection{Bayesian causal inference}
            \cite{samad-2015-perception-of-body-ownership-is-driven-by-bayesian-sensory-inference} modeled body ownership as Bayesian causal inference over visual, proprioceptive, and tactile signals. The model determines whether sensory observations are likely to arise from a common bodily cause or from independent causes, and it reproduced key properties of the rubber-hand illusion. This provides a computational account in which self-related sensory signals are integrated when a common cause is sufficiently probable and segregated otherwise. 


   \bibliographystyle{plainnat}
   \bibliography{/home/donkarlo/Dropbox/repo/nd_shared_research/bibliography/bibliography.bib}
\end{document}
